\documentclass{article}
\usepackage[utf8]{inputenc}
\usepackage{hyperref}
\usepackage{caption}
\usepackage{graphicx}
\usepackage{geometry}
\usepackage{xcolor}
\usepackage{fontspec}
\usepackage{listings}
\usepackage{float}



\setmainfont{Arial}
\geometry{a4paper, margin=1in}

\lstdefinestyle{JavaStyle}{
	language=Java,
	basicstyle=\ttfamily\small,
	keywordstyle=\color{blue}\bfseries,
	commentstyle=\color{gray}\itshape,
	stringstyle=\color{red!70!black},
	numbers=left,
	numberstyle=\tiny\color{gray},
	stepnumber=1,
	numbersep=8pt,
	showstringspaces=false,
	breaklines=true,
	breakatwhitespace=true,
	tabsize=4,
	frame=single,
	rulecolor=\color{black!25},
	captionpos=b,
	keepspaces=true,
	showtabs=false,
	showspaces=false
}

\lstset{style=JavaStyle}


\begin{document}

\section{Descripción}
% De que trata el proyecto

\section{Código}
\subsection{Archivos}
% Explicacion de las rutas de los archivos y que hace cada uno

\subsection{Manejo de hilos}
% Implementacion de hilos, sincronizacion, etc.
\subsection{Gestión de memoria}
% Explicacion de como se maneja la memoria, estructuras de datos, etc.

\subsection{Seguridad}
% Medidas de seguridad implementadas, manejo de errores, etc.
\subsubsection{}




\end{document}